\input{estilo}
% La lectura en el 11 de tabla-del-mvcc, como línea de tiempo: una línea por
% clave sobre el eje de timestamps (x = t − 8). La lectura toma, en cada clave,
% la última versión a su izquierda; la del 12 queda en el futuro y no se lee,
% pero no hace falta abortar.
\begin{document}
\begin{figura}\begin{tikzpicture}[x=0.9cm]
  \newcommand\val[2][]{\var{V}\textsubscript{\var{#2}}\ensuremath{#1}}
  \node[anchor=west] at (-0.4,1.2) {\var{x}};  \node[anchor=west] at (-0.4,0) {\var{y}};
  \draw[suave, line width=0.5pt] (0.2,1.2) -- (6.6,1.2);  \draw[suave, line width=0.5pt] (0.2,0) -- (6.6,0);
  \node[evento] (x) at (2,1.2) {};  \node[above left=1pt] at (x) {\val{x}};
  \node[evento] (y9) at (1,0) {};   \node[above left=1pt] at (y9) {\val{y}};
  \node[evento=suave] (y12) at (4,0) {};  \node[above right=1pt, text=suave] at (y12) {\val[']{y}};
  \node[etiq, text=suave, anchor=west, align=left] at (6.8,0.6) {varias versiones\\por clave};
  \draw[segundo, densely dashed] (3,-0.5) -- (3,1.8) node[above, text=segundo] {R/O @ 11};
  \draw[->, segundo] (3,1.05) -- (2.15,1.05);
  \draw[->, segundo] (3,-0.15) -- (1.15,-0.15);
  \node[etiq, text=suave, anchor=north west, align=left] at (4.15,-0.1) {futuro:\\no se lee};
  \tiempo[timestamp]{0.2}{9.0}{-1.0}
  \foreach \t in {9,10,11,12} {\hora{(\t-8,-1.0)}{\t}}
\end{tikzpicture}\end{figura}
\end{document}
